\documentclass[11pt]{article}
\usepackage[a4paper,margin=23mm]{geometry}
\usepackage{amsmath,amssymb,amsthm,booktabs,hyperref,xurl}
\hypersetup{hidelinks}
\setlength{\emergencystretch}{2em}
\newtheorem*{proposition}{Proposition}
\begin{document}
\begin{center}
{\Large One product identity removes a quadratic coordinate}\\[6pt]
\(\mathbb Z\mathbb Q\) Diophantine Project\qquad 23 November 2026
\end{center}

The corrected source stores \(R=\texttt{sr}=r\,s\,n^2\), together
with \(U=w\,n^2\), and has the three quadratic residuals
\[
 f_{\rm sq}=R_2-R^2,\qquad
 f_a=a-(U+1)R,\qquad f_p=p-2UR_2.
\]
The register \(R_2=\texttt{sr2}\) occurs only in \(f_{\rm sq}\)
and \(f_p\). Retain \(f_a\), delete \(R_2,f_{\rm sq}\), and replace
\(f_p\) by the quadratic residual
\[
                         g_p=p-2aR+2R^2.
\]

\begin{proposition}
This transformation is an exact equivalence on the entire rational
zero set after forgetting \(R_2\). Its inverse lift is unique and
polynomial. It saves one coordinate and one residual, independently
of any integer predicate or nonzero guard.
\end{proposition}
\begin{proof}
The identity
\[
                  g_p=f_p+2U f_{\rm sq}-2R f_a
\]
shows that every old solution gives a new solution. Conversely,
assign \(R_2=R^2\). Then \(f_{\rm sq}=0\), and the same identity
gives \(f_p=g_p+2R f_a=0\). The lift is forced by the deleted
square equation. No division is used; in particular \(R=0\),
\(U=0\), and \(a=0\) cause no exceptional branch. Every other
residual is copied unchanged after coordinate renumbering.
\end{proof}

Applying the transformation to the latest complete tables gives:
\begin{center}
\begin{tabular}{lrrrr}
\toprule
Interface & Old variables & New variables & Old rows & New rows\\
\midrule
General & 973 & 972 & 827 & 826\\
Specified & 1147 & 1146 & 997 & 996\\
Ordinary \(abc\) & 697 & 696 & 592 & 591\\
\bottomrule
\end{tabular}
\end{center}
The general and ordinary-\(abc\) predecessors are the complete
signed-power tables in \nolinkurl{undecidability/pell_signed_q_v98/};
the specified predecessor is
\nolinkurl{undecidability/pell_halftrace_v96/halftrace_m179_quadrics.json}.
The sums of squares still have degree exactly four, since retained
nonzero quadratic homogeneous parts cannot cancel in a real sum of
squares. All arithmetic-interface premises, instance ranges, integer
inputs, and the cited Jones coding theorem are unchanged. This is an
unconditional algebraic simplification of those systems; it introduces
no unconditional undecidability assertion.

The three full tables are \nolinkurl{product_general_quadrics.json},
\nolinkurl{product_specified_quadrics.json}, and
\nolinkurl{product_abc_quadrics.json}, each paired with a certificate
in the present directory. Run
\begin{verbatim}
python undecidability/pell_product_v99/build_product.py
python undecidability/pell_product_v99/verify_product.py
\end{verbatim}
The independent checker imports no builder. It verifies both polynomial
identities, the two old occurrences of \(R_2\), the exact replacement
row, every unchanged row, coefficient integrality, complete support,
degree, metadata and file hashes. The recovery conventions
\(\texttt{M\_recovery}=\texttt{sr}\) and
\(\texttt{a\_recovery}=(U+1)\texttt{sr}\) are retained; the obsolete
square-register description is preserved only as predecessor history.
\end{document}
